\documentclass[11pt,a4paper]{article}

% --- Packages ---
\usepackage[margin=.75in]{geometry}
\usepackage[utf8]{inputenc}
\usepackage[T1]{fontenc}
\usepackage{lmodern} % Fix for font not found error
\usepackage{titlesec}
\usepackage{enumitem}
\usepackage{hyperref}
\usepackage{xcolor}
\usepackage{fancyhdr}
\usepackage{comment}
\usepackage{cleveref}
\usepackage[bottom]{footmisc}
\usepackage{graphicx}
\usepackage{longtable}
\usepackage{booktabs}
\usepackage{array}
\usepackage{calc}
\usepackage{etoolbox}

% Standard Pandoc setup
\providecommand{\tightlist}{%
  \setlength{\itemsep}{0pt}\setlength{\parskip}{0pt}}

\crefformat{section}{\S#2#1#3}
\crefformat{subsection}{\S#2#1#3}
\crefformat{subsubsection}{\S#2#1#3}
\crefformat{figure}{#2Figure~#1#3}
\crefformat{footnote}{#2\footnotemark[#1]#3}

\linespread{0.9} % Riduce lo spazio tra le linee

% --- Colors ---
\definecolor{primary}{RGB}{0, 0, 0}
\definecolor{links}{HTML}{0000EE}
\definecolor{blue}{HTML}{26428B}

\hypersetup{
    colorlinks=true,
    linkcolor=links,
    urlcolor=links,
    citecolor=links,
}

% --- Section Formatting ---
\titleformat{\section}
  {\normalfont\Large\bfseries\color{blue}}{\thesection}{1em}{}
\titleformat{\subsection}
{\normalfont\large\bfseries\color{blue}}{\thesubsection}{1em}{}
\titleformat{\subsubsection}
{\normalfont\bfseries\color{blue}}{\thesubsubsection}{1em}{}

% --- Header and Footer ---
\pagestyle{fancy}
\fancyhead[L]{\small Giuseppe Capasso}
\fancyhead[R]{\small vapt}
\fancyfoot[C]{\thepage}
\renewcommand{\headrulewidth}{0.4pt}

% --- Custom Title Command ---
\newcommand{\proposaltitle}[1]{
    \begin{center}
        \vspace*{0.1cm}
        {\LARGE \bfseries \color{blue} #1} \\[1.5ex]
        {\large \textbf{Giuseppe Capasso}} \\[1ex]
                \small{
            \vspace{0.1cm}
            \textbf{Materia:} vapt\\
        }
            \end{center}
    \vspace{0.3cm}
    \hrule height 0.5pt
    \vspace{0.5cm}
}

\begin{document}

\proposaltitle{Andrea BIONDO}

Andrea BIONDO Università di Padova

Software Security 1 1

https://cybersecnatlab.it

License \& Disclaimer 2

License Information

Disclaimer

This presentation is licensed under the Creative Commons BY-NC License

To view a copy of the license, visit:
http://creativecommons.org/licenses/by-nc/3.0/legalcode

© CINI -- 2021

➢

We disclaim any warranties or representations as to the accuracy or
completeness of this material.

➢

Materials are provided ``as is'' without warranty of any kind, either
express or implied, including without limitation, warranties of
merchantability, fitness for a particular purpose, and non-infringement.

➢

Under no circumstances shall we be liable for any loss, damage,
liability or expense incurred or suffered which is claimed to have
resulted from use of this material.

Rel. 22.03.2021

Obiettivi 3

Comprensione della memoria a basso livello ➢ Conoscenza di base dei
buffer overflows ➢ Conoscenze di base di reverse engineering ➢

© CINI -- 2021

Rel. 22.03.2021

Prerequisiti 4

➢

Sistemi di numerazione ➢ Esadecimale

➢

Conoscenza di base del linguaggio C ➢ Tipi primitivi e strutturati,
puntatori ➢ Funzioni di I/O

© CINI -- 2021

Rel. 22.03.2021

Argomenti 5

Spazio di memoria ➢ Reverse engineering ➢ Buffer overflows ➢

© CINI -- 2021

Rel. 22.03.2021

Argomenti 6

Spazio di memoria ➢ Reverse engineering ➢ Buffer overflows ➢

© CINI -- 2021

Rel. 22.03.2021

Cos'è la memoria? 7

Per un programmatore potrebbe essere un insieme di variabili tipate ➢
Per un ingegnere elettronico potrebbe essere un insieme di celle
bistabili ➢ Dobbiamo scegliere un livello di astrazione ➢

© CINI -- 2021

Rel. 22.03.2021

Astrazioni di memoria 8

Dati tipati (variabili)

Visione interpretata dei byte

Linguaggio di programmazione Sequenza di byte indirizzabili Spazio
indipendente per-processo Solo alcune aree sono mappate

Memoria virtuale Sistema operativo

Memoria fisica © CINI -- 2021

Sequenza di byte indirizzabili

Rel. 22.03.2021

Memoria virtuale 9

Spazio virtuale di dimensione fissata (4GB / 256TB) ➢ Mapping fra aree
di memoria virtuale e fisica ➢ Flag di protezione degli accessi ➢

➢ Read, write, execute

© CINI -- 2021

Rel. 22.03.2021

Spazio virtuale Linux userspace 10

Indirizzi

Binario Text

Codice

Data

Dati globali inizializzati

BSS

Dati globali azzerati

Heap

Allocazioni dinamiche

Librerie

Binari delle librerie dinamiche

Stack

Variabili locali, record di attivazione

© CINI -- 2021

Rel. 22.03.2021

Spazio virtuale Linux userspace 11

© CINI -- 2021

Rel. 22.03.2021

Rappresentazione di interi 12

➢

unsigned int: intero a 32 bit senza segno ➢ Numero in {[}0,
4294967295{]}

➢

Esempio: valore 100.000.000 ➢ Esadecimale: 0x05F5E100 +0

Big-endian

+1

0x05 +0

Little-endian

+2

0xF5 +1

0x00

+3

0xE1 +2

0xE1

0x00 +3

0xF5

© CINI -- 2021

0x05 Rel. 22.03.2021

x86 usa LE

Rappresentazione di tipi C (x86) 13

➢

Interi little-endian ➢ Interi con segno rappresentati secondo la
notazione

complemento a due ➢ char, int, short, long, enum: interi a varie
lunghezze

float e double: IEEE 754 ➢ I puntatori sono interi unsigned ➢

➢ Il loro valore è l'indirizzo puntato ➢ 32/64 bit (per indirizzare
intero spazio virtuale)

© CINI -- 2021

Rel. 22.03.2021

Rappresentazione di tipi C (x86) 14

➢

Array ➢ Elementi disposti sequenzialmente ➢ {[}i{]} @ base array + i *
size elemento

➢

Strutture ➢ Campi disposti sequenzialmente in ordine di dichiarazione ➢
Campo @ base struct + somma size campi precedenti ➢ (Il compilatore
potrebbe introdurre del padding) © CINI -- 2021

Rel. 22.03.2021

Corruzione di memoria 15

Modificare la memoria di un processo in un modo diverso da quello
previsto dal programmatore ➢ Controllare la memoria = controllare il
processo! ➢

© CINI -- 2021

Rel. 22.03.2021

Corruzione di memoria in the wild 16

➢

Malware ➢ Morris worm (1988!), Blaster, Sasser, Conficker, \ldots{} ➢
Più recentemente, StuxNet e WannaCry

➢

Servizi remoti e applicazioni utente ➢ Attacchi su server e browser

➢

Sbloccaggio di dispositivi ➢ Root Android, jailbreak iOS, console per
videogiochi

© CINI -- 2021

Rel. 22.03.2021

Argomenti 17

Spazio di memoria ➢ Reverse engineering ➢ Buffer overflows ➢

© CINI -- 2021

Rel. 22.03.2021

Reverse engineering 18

``Analizzare un sistema per creare rappresentazioni ad alto livello di
astrazione'' (Chikofsky, Cross 1990)

Prodotto finito (eseguibile compilato)

© CINI -- 2021

Informazioni progettuali (architettura, sorgente)

Rel. 22.03.2021

Reverse engineering 19

➢

Molte ragioni per fare reversing ➢ Recupero di codice sorgente ➢
Documentazione mancante o insufficiente ➢ Analisi di prodotti
concorrenti ➢ ``Aprire'' piattaforme proprietarie ➢ Auditing di
sicurezza ➢ Curiosità © CINI -- 2021

Rel. 22.03.2021

La vita di un programma 20

Sorgente

Perdita di informazioni ad alto livello

Compilatore

Analisi dinamica

File oggetto

Processo

Linker

Eseguibile

Lib. statiche

Loader Linker dinamico Lib. dinamiche

Analisi statica © CINI -- 2021

Rel. 22.03.2021

Eseguibili 21

➢

Molti formati, a seconda dell'OS ➢ PE (Windows), Mach-O (MacOS, iOS),
ELF (*nix), \ldots{}

➢

Divisi in varie sezioni/segmenti mappabili ➢ Cioè che diventano aree di
memoria virtuale a runtime

© CINI -- 2021

Rel. 22.03.2021

Gli strumenti 22

➢

Analisi statica ➢ Disassemblatori, decompilatori (e.g., Ghidra) ➢
Avanzati: interpretazione astratta, esecuzione simbolica, \ldots{}

➢

Analisi dinamica ➢ Debugger (e.g., GDB) ➢ Avanzati: tracer,
instrumentazione dinamica, \ldots{}

© CINI -- 2021

Rel. 22.03.2021

Assembly x86\_64 23

La CPU ha una piccola memoria locale composta da registri ➢ Ogni
istruzione assembly ha degli operandi ➢

➢ Registri: rax, ebx, r13d, \ldots{} ➢ Memoria: {[}rax+4{]} ➢

Notazione Intel: , © CINI -- 2021

Rel. 22.03.2021

Registri x86\_64 24

Estesi da x86: r\{a,b,c,d\}x ➢ Program counter: rip ➢ Gestione stack ➢

➢ rsp (stack pointer) ➢ rbp (frame pointer) ➢

Generici: r8-r15 © CINI -- 2021

Rel. 22.03.2021

Registri x86\_64 25

RAX

EAX

AX

AH

AL

8

8

16 32 64

© CINI -- 2021

Rel. 22.03.2021

Alcune istruzioni di base 26

MOV , ➢ PUSH / POP ➢ ADD/SUB , ➢ CALL / RET ➢

© CINI -- 2021

Rel. 22.03.2021

Salti condizionali 27

➢

CMP , ➢ Confronta due valori e imposta delle flag

➢

J ➢ Salta a se le flag matchano

➢

Salta se rax != 15: ➢ CMP rax, 15

➢ JNE \ldots{}

© CINI -- 2021

Rel. 22.03.2021

Reversing statico con Ghidra 28

© CINI -- 2021

Rel. 22.03.2021

Argomenti 29

Spazio di memoria ➢ Reverse engineering ➢ Buffer overflows ➢

© CINI -- 2021

Rel. 22.03.2021

Buffer overflows 30

char name{[}100{]}; printf(``Come ti chiami?''); scanf(``\%s'', name);
printf(``Ciao, \%s!\n'', name);

© CINI -- 2021

Rel. 22.03.2021

Buffer overflows 31

➢

Cosa succede se l'utente ha un nome più lungo di 100 caratteri? ➢ O se è
malevolo e vuole dare più di 100 caratteri in input\ldots{}

➢

scanf (in questo caso) non controlla i bound ➢ Continuerà a scrivere
caratteri oltre la fine di name,

sovrascrivendo la memoria che lo segue

© CINI -- 2021

Rel. 22.03.2021

Buffer overflows 32

➢

Si ha un buffer overflow quando il programma scrive oltre la fine di un
buffer ➢ Perché i dati sono più lunghi del buffer

➢

Classica vulnerabilità di corruzione della memoria ➢ Scriviamo dati
dall'attaccante (input utente) in posizioni di

memoria che il programmatore non aveva previsto potessero essere
modificate (oltre la fine del buffer) © CINI -- 2021

Rel. 22.03.2021

Conseguenze del buffer overflow 33

Le garanzie di correttezza cadono completamente ➢ Se corrompiamo
abilmente la memoria, possiamo ottenere il controllo completo del
processo ➢

➢ Arbitrary code execution

© CINI -- 2021

Rel. 22.03.2021

Accessi out-of-bounds 34

struct \{ int a{[}2{]}; int b{[}3{]}; \};

a{[}0{]} a{[}1{]} b{[}0{]} b{[}1{]} b{[}2{]}

Questo è anche a{[}2{]}!

a{[}2{]} = 42; printf(``\%d\n'', b{[}0{]}); /* stampa 42 */ © CINI --
2021

Rel. 22.03.2021

Accessi out-of-bounds 35

Nessun bound check! Scritture OOB

© CINI -- 2021

Rel. 22.03.2021

Accessi out-of-bounds 36

Nessun bound check!

Scritture OOB

© CINI -- 2021

Rel. 22.03.2021

Pattern pericolosi 37

➢

Senza bound checking ➢ gets, scanf \%s, sprint, strcpy

➢

Bound checking usato impropriamente ➢ fgets, (f)read, snprintf, memcpy,
strncpy, \ldots{}

➢

Manipolazioni manuali ➢ Loop di copia, accessi ad array, \ldots{}

© CINI -- 2021

Rel. 22.03.2021

Un primo overflow 38

Nessun bound check!

Overflow dei caratteri 89

© CINI -- 2021

Rel. 22.03.2021

Andrea BIONDO Università di Padova

Software Security 1 39

https://cybersecnatlab.it



\end{document}
